\documentclass[11pt,a4paper]{article}

\usepackage{graphicx}   % only required graphics package for DVI and PDF
\usepackage{natbib}    % normally combine with: \bibliographystyle{apacite}
%\usepackage[numbers]{natbib}    % numbers, normally combine with: \bibliographystyle{aip}
\bibliographystyle{apalike}

%\setlength{\bibsep}{0.1pt}      % shorten separation in references list
\usepackage[bf]{caption}   % allows for caption text options etc
\usepackage{booktabs}   % cool tables
% \usepackage{ctable}     % allow flexible captions and footnotes for tables

% include all math symbol packages by default
\usepackage{amsmath}
\usepackage{amssymb}
\usepackage{latexsym}

\usepackage{hyperref}
\hypersetup{
	colorlinks=true,
	linkcolor=blue,
}

% do spacing
% vital for accurate page layouts
\usepackage[margin=3cm]{geometry}

\usepackage{microtype} % package for tidying up text and font issues; does a beautiful job of stopping text overlapping the next column, or running into the whitespace

%%% option for double-spacing
%\renewcommand{\baselinestretch}{1.66}

% page-numbers at top
\pagestyle{myheadings}

\begin{document}
%%%%%%%%%%% TITLE SECTION %%%%%%%%%%%%%%%%%%%%%%%%%%%%%%%
\noindent \Large{Analysis of stochastic choice in foraging} \\

\vspace{0.25cm}

\normalsize
\noindent Mark D. Humphries \\

\noindent Version 1.0 \\
\noindent Started: 28 July 2023 \\

% \vspace{0.5cm}

\noindent {\bf Address:} School of Psychology, University of Nottingham \\

\noindent {\bf Corresponding author:} mark.humphries@nottingham.ac.uk \\

\begin{abstract}
Optimal models of patch foraging assume a deterministic choice rule. We analyse here a stochastic choice model of patch foraging, based on exploit-explore action selection models from reinforcement learning.
\end{abstract}

%%%%%%%%%% MAIN BODY %%%%%%%%%%%%%%%%%%%%%%%%%%%%%%%%%%%%%
\section{Introduction}
We introduce here a model of patch foraging that uses stochastic choice for its stay-or-leave decision. In particular, we propose a model that has parametric control over the exploration-exploitation tradeoff. We analyse its expected leaving times (or, equivalently, its dwell time in each patch), and the variance in leaving time as a function of these parameters.

We assume here discrete time-steps (or ``trials'') throughout, as this allows a direct comparison with the simulated model(s) used to fit participant data. It also makes the maths easier. The translation to continuous time is considered in Section \ref{section:continuous}.  \\

\noindent \textbf{Notation for the models} \\
$n$: the integer time-step (or trial number) \\
$t(n)$: the clock time of time-step $n$ \\
$r(n)$: the amount of reward in the patch at time-step $n$ \\
$r_0$: the initial amount of reward in the patch before harvesting \\
$\alpha$: the decay rate of $r(n)$ \\
$R$: the average or ``background'' reward rate in the environment \\ 
$\beta$: inverse temperature parameter of the softmax function (higher values mean more exploitation -- lower probability of leaving)
$\epsilon$: fixed probability of leaving per time-step

\section{The stochastic choice model} 

We define models for the probability of leaving a patch on the $n$th time-step given the current reward rate $r(n)$. These models are derived from their reinforcement-learning equivalents.

The softmax choice model is:
\begin{equation} \label{eq:softmax}
p(\mathrm{leave}|n) = \frac{1}{1+\exp(\beta r(n))}.
\end{equation}

The $\epsilon$-greedy model is
\begin{equation}\label{eq:epsilon}
p(\mathrm{leave}|n) = \epsilon,
\end{equation}
a fixed probability of leaving per time-step.

\section{Reward functions}
We consider here two standard decaying reward functions.

Exponential decay is given by
\begin{equation} \label{eq:expdecay}
r(n) = r_0\exp(-\alpha n).
\end{equation}

Linear decay is given by 
\begin{equation} \label{eq:lineardecay}
r(n) = \begin{cases} 
r_0 -\alpha n & \mathrm{r_0 \geq \alpha n} \\
0 & \mathrm{otherwise}
\end{cases} 
\end{equation}

\section{MVT predictions for leaving time}
The key metric for foraging is the leaving time for a patch i.e. how long an agent will stay in the patch before leaving. Optimal foraging theory (MVT) says that agents should leave as soon as the reward rate in the patch $r(n)$ is equal to the long-run average rate $R$ for the environment. MVT predicts independent patch and environment effects on leaving time. It predicts leaving times will be longer for richer patches (the patch effect). Conversely, it predicts leaving times for a given patch will be longer in a poorer environment, one with a lower average reward rate. For a given patch we can calculate the predicted MVT leaving time assuming $R$ is known.

[THESE NEED DOUBLE-CHECKING]

For the exponentially-decaying reward function we set $R = r_0\exp(-\alpha n)$ and solve for $n$
\begin{equation}
n = -\log\left(\frac{R}{r_0}\right)\frac{1}{\alpha}.
\end{equation}

For the linearly-decaying reward function we set $R = r_0 -\alpha n$ and solve for $n$
\begin{equation}
n = \frac{r_0 - r}{\alpha},
\end{equation}
where we can ignore the second case in Eq. \ref{eq:lineardecay} as we can guarantee that $R \geq 0$. 

Whatever the reward function, if $n$ is a discretisation of continuous time, then the expected time is $t(n)$.


\section{How to compute expectation and variance for leaving times}
We want our models to replicate both patch and environment effects. We also would like them to explain overharvesting: that agents stay longer in patches than predicted by MVT. Consequently, we want to calculate the expected leaving time for each model. But given the choice is stochastic we also want the predicted variation around that expected time. We do this as follows.

The models give us the probability of leaving $p(\mathrm{leave}|n)$ on a given time-step $n$. We want to know the probability $p(\mathrm{leave} = n)$  that time-step $n$ is the first trial at which leaving occurs. Each time-step $n$ has only two possible outcomes (stay or leave) and is thus a Bernoulli trial; the sequence of time-steps $n$ is a Bernoulli process. The probability $p(\mathrm{leave} = n)$ is thus
\begin{equation}\label{eq:p_leave_now}
p(\mathrm{leave} = n) = p(\mathrm{leave}|n) \prod_{i=1}^{n-1} 1-p(\mathrm{leave}|i),
\end{equation}  
where the first term is the probability of leaving on this time-step and the second term is the cumulative probability of not having left on all the previous time-steps.

The tasks we model can have discrete trials or states in which a stay-or-leave decision is made after every harvest in a patch. We can thus compute exact values of expectation and variance by plugging Eq. \ref{eq:p_leave_now} into the standard equations for expectation ($\sum xp(x)$) and variance ($\sum (x-E)^2p(x)$) of a random variable ($x$). Here the random variable is $n$, the number of elapsed trials.

The expected leaving time is therefore
\begin{equation}\label{eq:E_leave_discrete}
\mathrm{E(leave)}  = \sum_{n=1}^{\infty} n p(\mathrm{leave} = n).
\end{equation} 

Similarly, the variance in leaving time is therefore
\begin{equation}\label{eq:VAR_leave_discrete}
\mathrm{VAR(leave)}  = \sum_{n=1}^{\infty} \left(n - \mathrm{E(leave)})\right)^2 p(\mathrm{leave} = n).
\end{equation} 

In special cases these equations have exact solutions. Otherwise, to evaluate these numerically, we replace $\infty$ with some large number $n_{\textrm{max}}$ as the upper limit of time-steps over which to evaluate the expectation. We need to choose that limit such that $n_{\textrm{max}} \gg \mathrm{E(leave)}$. 

\subsection{Continuous time} \label{section:continuous}
For tasks with continuous time in a patch we can treat $n$ as discrete time-steps approximating the continuous time, to match how such tasks are handled in the simulated models. Each time-step is $\delta  = t(n) - t(n-1)$; equivalently: $t(n) = n \delta$.

The discrete probabilities defined above are for each trial, which have unit time (i.e. $t(n) = n$). So if we change the time-step $\delta$ from unit time then we also need to scale the probability of leaving $p(\mathrm{leave}|n)$ to get the same probability per unit time. Namely:
\begin{equation}\label{eq:p_leave_cont}
p(\mathrm{leave}|n)_\delta  = p(\mathrm{leave}|n) \delta.
\end{equation} 

We plug $p(\mathrm{leave}|n)_\delta$ into equation \ref{eq:p_leave_now} as before, obtaining $p(\mathrm{leave} = n)_\delta$. We then compute the expectation and variance using the clock time $t(n)$ of each time-step
\begin{equation}\label{eq:E_leave_cont}
\mathrm{E(leave)}  = \sum_{n=1}^{\infty} t(n)p(\mathrm{leave} = n)_\delta.
\end{equation} 

\begin{equation}\label{eq:VAR_leave_cont}
\mathrm{VAR(leave)}  = \sum_{n=1}^{\infty} \left(t(n) - \mathrm{E(leave)})\right)^2 p(\mathrm{leave} = n)_\delta.
\end{equation} 
Numerical calculation of these again uses some upper limit $n_{\textrm{max}}$ in place of $\infty$. We need to choose that limit such that $t(n_{\textrm{max}}) \gg \mathrm{E(leave)}$. 

Numerical simulations confirm that $\mathrm{E(leave)}$ and $\mathrm{VAR(leave)}$ converge on stable values as $\delta \rightarrow 0$. Nonetheless it would be good to handle continuous time directly, rather than as the approximations above. The limiting case of the taking $\delta t \rightarrow 0$ for a Bernoulli process is the Poisson process, which we will come back to at a later date. 

\section{Choice behaviour for the $\epsilon$-greedy model}
This model has a fixed probability of leaving per time-step. The probability of leaving (Eq. \ref{eq:p_leave_now}) is thus independent of the reward function $r(n)$  
\begin{align}
p(\mathrm{leave} = n) & = \epsilon \prod_{i=1}^{n-1} 1-\epsilon, \nonumber \\
& = \epsilon(1-\epsilon)^{n-1}. \label{eq:epsilon_p_leave}
\end{align}  

Consequently, the expectation and variation in leaving time are also independent of the reward function $r(n)$. They are found by plugging equation (\ref{eq:epsilon_p_leave}) into equations (\ref{eq:E_leave_discrete}) and (\ref{eq:VAR_leave_discrete}) for discrete trials or equations (\ref{eq:E_leave_cont}) and (\ref{eq:VAR_leave_cont}) for continuous time.

If we have discrete trials $n$ then equation (\ref{eq:epsilon_p_leave}) is a geometric distribution with exact solutions for its expectation and variance (see e.g. \href{https://en.wikipedia.org/wiki/Geometric\_distribution}{Wikipedia})
\begin{align*}\label{eq:geometric_distribution}
\mathrm{E(leave)} & = \frac{1}{\epsilon} \\
\mathrm{VAR(leave)} & = \frac{1-\epsilon}{\epsilon^2}.
\end{align*}
It is thus obvious that:
\begin{itemize}
	\item There is a monotonic relationship between $\epsilon$ and leaving time: increasing $\epsilon$ shortens the leaving time (as expected)
	\item There is a monotonic relationship between $\epsilon$ and the variation in leaving time: increasing $\epsilon$ reduces the variance. 
	\item There is no patch effect as changing $r_0$ will not affect $\mathrm{E(leave)}$
	\item The environment effect can be captured by changing $\epsilon$ 
\end{itemize}

\section{Choice behaviour for the softmax model}

\subsection{Exponentially decaying rewards}
Putting the exponential reward decay (Eq. \ref{eq:expdecay}) into the softmax model gives:
\begin{equation} \label{eq:softmax_p_leave}
p(\mathrm{leave}|n) = \frac{1}{1+\exp(\beta r_0\exp(-\alpha n))}.
\end{equation}
It is immediately clear that $\beta$ and $r_0$ together form single parameter $B = \beta r_0$ that sets the explore-exploit tradeoff in foraging. As $\beta$ and $r_0$ can be changed independently, this could allow the environment and patch effect. The environment effect between patches of the same type could thus be achieved by changing $\beta$ with $r_0$ held constant. The patch effect in a given environment is thus a change in $r_0$ bewteen patches with $\beta$ held constant. 


% \bibliography{AnalysisRefs} % bibliography file

\end{document}
